\section{Method and Protocol}
\label{sec:method}

\subsection{Backbone, data and baselines}
We fine-tune the non-distilled \deit checkpoint (\texttt{deit-small-patch16-224}: 12 blocks, width 384, 6 heads, 22M parameters)~\cite{touvron2021deit} independently on three MedMNIST~v2 tasks at $224\times224$~\cite{yang2023medmnist}: \pathmnist (9 colorectal tissue classes; 90k/10k/7.2k train/val/test), \dermamnist (7 skin lesion classes; 7.0k/1.0k/2.0k) and \pneumnist (binary chest X-ray; 4.7k/0.5k/0.6k). Official splits are kept; archives are MD5-verified. Training uses random-resized crops (scale 0.8--1.0) and horizontal flips, AdamW~\cite{loshchilov2019adamw} with learning rate $5\times10^{-5}$, weight decay 0.05, batch 32 and a cosine schedule. Epoch counts (5, 10, 5) were fixed from single-seed pilots on validation. Three seeds (0, 1, 2) are trained per dataset and the best-validation checkpoint is kept; test data never informs any choice. A loading check rejects any missing or mismatched backbone weight, since an early run that silently dropped the pretrained weights produced misleading baselines.

\subsection{Token reduction adapters}
All adapters act on the frozen fine-tuned model and physically shorten the token sequence after a complete transformer block (block 4 unless stated), keeping the class token and original positional embeddings. Budgets retain $n\in\{147, 98, 69, 49, 20\}$ of 196 image tokens, i.e.\ 25--90\% reduction. Because earlier blocks process all tokens, $n$ measures sequence length downstream of the reduction point, not total compute or discarded image content; we report measured FLOPs separately.

\textbf{Random deletion} ranks tokens with i.i.d.\ uniform scores; three mask seeds separate mask variance from training variance. \textbf{Attention deletion} keeps the top-$n$ tokens by head-averaged class-query attention from the preceding block; the cost of recomputing that row is included in profiling. \textbf{EViT-style fusion}~\cite{liang2022evit} keeps the top-$(n{-}1)$ attention-ranked tokens and replaces the rest with one attention-weighted summary that counts toward the budget. \textbf{ToMe-style merging}~\cite{bolya2023tome} repeatedly splits image tokens into alternating sets, matches by cosine similarity of block-output features, and merges the most similar pairs until $n$ remain; merged tokens carry mass, which biases later attention logits by $\log$ mass. Fusion happens after a full block rather than between attention and MLP, and merging repeats at one depth rather than at every layer, so both are controlled adaptations rather than reproductions. A no-op adapter is tested bit-exact against the original model.

\subsection{Fine-tuning and recovery methods}
\label{sec:finetune}
We compare four ways of adapting the backbone, each with a learning rate chosen on a seed-0 validation pilot: \textbf{full} ($5\times10^{-5}$), \textbf{partial} (last 4 blocks and final norm, $3\times10^{-4}$), \textbf{linear} (head only, $3\times10^{-3}$) and \textbf{LoRA}~\cite{hu2022lora} on the query and value projections ($r{=}8$, $\alpha{=}16$, dropout 0.1, $1\times10^{-3}$) via PEFT~\cite{peft}. The head is always trained. LoRA weights are merged on export so every checkpoint is a plain ViT state dict usable by the same evaluation and pruning code. For every run we record trainable parameters, optimisation-only wall-clock time, peak GPU memory, and validation accuracy four times per epoch to estimate time-to-target.

\textbf{Recovery} re-trains a pruned model for two epochs with the reduction active, starting from the fine-tuned checkpoint, at a budget one step more aggressive than the validation-selected one (floor 20 tokens). A full-token control receives the same extra training from the same checkpoint, so the recovery comparison is matched in training budget. We run recovery with full-parameter updates ($10^{-5}$) and with LoRA-only updates ($2\times10^{-4}$), and also test LoRA recovery from a LoRA-fine-tuned start, i.e.\ a parameter-efficient pipeline end to end.

\subsection{Selection, locking and confirmation}
\label{sec:protocol}
The protocol, including tolerances, seeds and the interval procedure, is written to a JSON file before the test stage. The primary tolerance is $\delta{=}0.5$\,\pp of accuracy, with 1\,\pp as a secondary margin; these are research thresholds, not clinical criteria. On validation, a method's selected budget is the most aggressive $n$ whose accuracy is within $\delta$ of the full model for \emph{every} training seed (and every mask seed for random). Ties across methods prefer more reduction, then mean accuracy, then latency. For the selected method we ablate reduction depth (after blocks 1, 4, 8), a progressive schedule (blocks 4, 7, 10 with geometrically decreasing counts) against a single-stage control matched in FLOPs, and input-stage removal before block 1 (random and uniform grid) with the full patch-embedding cost still counted.

Before any test inference, the selected policies, their specification hashes, the protocol hash and all 27 checkpoint hashes are committed to a manifest. For each fixed (reference, reduced) checkpoint pair the test accuracy difference is the fraction of images improved minus the fraction harmed. A paired image bootstrap (10{,}000 resamples) gives a descriptive 95\% interval~\cite{efron1994bootstrap}. The confirmatory statistic is a conservative lower bound: the Clopper--Pearson~\cite{clopper1934} lower bound on the improvement probability minus the upper bound on the harm probability, each at $\alpha/6$ with $\alpha{=}0.05$, which is simultaneous across the two tails and three seeds. Non-inferiority is supported only if all three seeds' bounds exceed $-\delta$~\cite{walker2011noninferiority}. Image rows are the resampling unit because MedMNIST does not expose patient identifiers.

\subsection{Efficiency measurement}
FLOPs count supported PyTorch operators at two per multiply-add and include selection, fusion and merging overhead; the image-token count entering every block is logged. Latency (batch 1) and throughput (batch 32) are medians of 30 timed forwards after 10 warm-ups on real held-out batches with device synchronisation, in float32 on an RTX PRO 6000 (Blackwell) GPU with PyTorch 2.11 and Transformers 5.20. Peak allocated memory is recorded. FLOP reduction and measured speed-up are always reported separately because they disagree.
